\documentclass{beamer}

\usepackage{../minimalbeamer}

\setbeamertemplate{footline}{
    \hspace{5mm}
    \hfill
    \insertframenumber{} / \inserttotalframenumber
    \hspace{8mm}
    \vspace{3mm}
}

\title{온라인 마트}
\author{}
\date{\number\year년 \number\month월 \number\day일}

\newcommand{\koreandate}{%
  \number\year년 \number\month월 \number\day일%
}
\begin{document}

\frame{\titlepage}

\begin{frame}{온라인마트}

\footnotesize

상품을 판매하는 온라인 마트가 있다. 상품은 ID로 구별되고 품목, 제조사, 가격 정보를 가지고 있다. 품목과 제조사는 1부터 5까지의 정수로 주어진다.

온라인 마트에서 상품을 판매할 수 있고 상품의 판매를 종료할 수 있다. 또한, 같은 품목과 제조사를 가진 상품들에 대해서 가격을 특정 금액만큼 할인할 수 있다.

예로 품목이 1이고 제조사가 2인 상품에서 대해서 5만큼 할인한다면 품목이 1이고 제조사가 2인 모든 상품의 가격이 5만큼 감소한다.

만약, 할인되어 가격이 0 또는 음수가 된 상품은 더 이상 판매가 의미 없는 상품이 되어 판매가 종료된다.

또한, 사용자에게 특정 조건을 만족하는 상품들에 대해서 가격이 낮은 순서로 최대 5개 상품을 검색하여 보여준다.

조건은 전체 상품, 특정 품목, 특정 제조사로 3가지 방식으로 줄 수 있다. 가격이 같은 경우 상품 ID가 더 낮은 것이 우선이다.

\end{frame}

\begin{frame}{\texttt{sell(mID, mCategory, mCompany, mPrice)}}
\textbf{상품 판매 시작}

\begin{itemize}
    \item ID가 $1 \le \text{mID} \le 10^9$인 상품 판매를 시작한다.
    \item 품목: $1 \le \text{mCategory} \le 5$
    \item 제조사: $1 \le \text{mCompany} \le 5$
    \item 가격: $1 \le \text{mPrice} \le 1\ 000\ 000$
    \item 동일한 \texttt{(mCategory, mCompany)}를 가진 판매 중 상품 개수를 반환한다.
    \item \texttt{mID}는 이전에 사용된 적 없는 값이 보장된다.
    \item 호출 횟수 최대 50,000회   
\end{itemize}
\end{frame}

\begin{frame}{closeSale(mID)}
\textbf{상품 판매 종료}

\begin{itemize}
    \item ID가 \texttt{mID}인 상품의 판매를 종료한다.
    \item 판매 종료 시 해당 상품의 가격을 반환한다.
    \item 다음 경우 \texttt{-1} 반환:
    \begin{itemize}
        \item 존재하지 않는 상품
        \item 이미 판매 종료된 상품
    \end{itemize}
    \item 호출 횟수 최대 5,000회
\end{itemize}
\end{frame}

\begin{frame}{discount(mCategory, mCompany, mAmount)}
\textbf{상품 할인}

\begin{itemize}
    \item 품목이 \texttt{mCategory}이고 제조사가 \texttt{mCompany}인 모든 상품 가격을 $1 \le \text{mAmount} \le 1\ 000\ 000$ 만큼 감소시킨다.
    \item 할인 후 가격이 0 이하가 되면 해당 상품은 판매 종료된다.
    \item 할인 후 동일한 \texttt{(mCategory, mCompany)}를 가진 판매 중 상품 개수를 반환한다.
    \item 호출 횟수 최대 10,000회
\end{itemize}
\end{frame}

\begin{frame}{show(mHow, mCode)}
\textbf{상품 조회}

\begin{itemize}
    \item 조건을 만족하는 판매 중 상품 중 가격이 낮은 순으로 최대 5개 조회. 가격이 같으면 ID가 작은 상품이 우선
    \item \texttt{mHow = 0}: 전체 상품 조회
    \item \texttt{mHow = 1}: 품목이 \texttt{mCode}인 상품 조회
    \item \texttt{mHow = 2}: 제조사가 \texttt{mCode}인 상품 조회
    \item 결과:
    \item 반환 개수 → \texttt{RESULT.cnt}
    \item i번째 상품 ID → \texttt{RESULT.IDs[i-1]}
    \item 조건 만족 상품이 없으면 \texttt{RESULT.cnt = 0}
    \item 호출 횟수 최대 1,000회
\end{itemize}
\end{frame}

\begin{frame}{관찰}
\begin{itemize}
    \item \texttt{sell()}을 50,000번 호출한 후에 \texttt{distcount()}를 10,000번 호출하는 경우가 존재하므로 단순 구현을 불가능합니다.
    \item (\texttt{mCategory}, \texttt{mCompany}) 범위가 작습니다.
    \item (\texttt{mCategory}, \texttt{mCompany}) 단위로만 업데이트가 일어납니다.
    \item \texttt{distcount()}는 언제나 가격을 낮추는 방향으로 업데이트가 일어납니다.
    \item \texttt{show()}는 같은 기준으로 가격이 낮은 순으로 최대 5개 상품을 반환합니다.
\end{itemize}
\end{frame}

\begin{frame}{풀이}
\begin{itemize}
    \item (\texttt{mCategory}, \texttt{mCompany}) 단위로 상품을 저장하고, 별도로 누적된 할인 금액을 저장합니다.
    \item 그런데, 가격이 0이하가 된 상품을 제거해야 하므로 최솟값을 빠르게 찾는 자료구조를 사용합니다.
    \item 새로 상품을 추가할때마다 기존 누적 할인 금액에 영향을 받지 않도록 상품의 가격에 할인 금액을 더해서 저장합니다.
    \item 임의의 상품을 제거해야하는 경우가 있기 때문에 최솟값을 빠르게 찾으면서 원소 제거가 가능한 treeset을 사용합니다.
    \item treeset 오버헤드로 인해 시간 초과가 발생하는 경우, 힙과 lazy deletion을 사용해서 뽑을때마다 이미 삭제된 원소인지 확인합니다.
\end{itemize}
\end{frame}

\begin{frame}[fragile, allowframebreaks]{김시온 소스 코드}
\inputminted{cpp}{../src/온라인마트/김시온/solution.cpp}
\end{frame}

\begin{frame}[fragile, allowframebreaks]{김지훈 소스 코드}
\inputminted{java}{../src/온라인마트/김지훈/UserSolution.java}
\end{frame}

\begin{frame}[fragile, allowframebreaks]{서민종 소스 코드}
\inputminted{java}{../src/온라인마트/서민종/UserSolution.java}
\end{frame}

\begin{frame}[fragile, allowframebreaks]{원찬혁 소스 코드}
\inputminted{java}{../src/온라인마트/원찬혁/UserSolution.java}
\end{frame}

\begin{frame}[fragile, allowframebreaks]{정의찬1 소스 코드}
\inputminted{java}{../src/온라인마트/정의찬/온라인마트v1.java}
\end{frame}

\begin{frame}[fragile, allowframebreaks]{정의찬2 소스 코드}
\inputminted{java}{../src/온라인마트/정의찬/온라인마트v2.java}
\end{frame}

\end{document}

